\section{The period Gram identity on every supplied original setting}
\label{sec:period-gram-v57}

The purpose of this chapter is to identify the actual Poincar\'e periods
with the original sum of orbital operators.  The identity is proved for
every supplied original setting and every pair of words, under exactly
the analytic premises of \texttt{KExpansionC}.  No arithmetic gap
inequality is added.  This is possible because that statement already
supplies the folding, square-integrability, orbital-integrability and
convergence hypotheses.  We prove the intervening measure decompositions
and finite-character identities; we do not replace them by another
period-matching assumption.

Brox's \emph{actual Hilbert-valued period Gram expansion}, at revision
\texttt{fd38313e}, suggested the useful
organization into support cells, rational returns and an antilinear
column map.  Its installed arithmetic gap is replaced here by the
literal universal analytic premises.  In particular, a summable series
of integrated operators is never used as a substitute for the absolute
integrability needed to unfold before integration.  That integrability
is proved separately from the supplied folding data.

\subsection{The fixed objects and the exact endpoint}

Fix an original datum \(D\), an arbitrary supplied
\(X:\texttt{HidaSetting}(D)\), and an arbitrary
\(Y:\texttt{Full}(X)\).  Write \(E/F\) for its CM extension,
\(G\) for its rational unitary group, and
\[
 G_{\mathbb A}=G_\infty\times G_f,\qquad
 \Gamma=G(F),\qquad d\mu_{\mathbb A}=d\mu_\infty\,d\mu_f.
\]
These are the original groups and measures, transported through the
supplied coefficient coordinates.  The rational embedding is the
original \texttt{RatLink}; we suppress that embedding in formulas.
Write \(r\) for the length of the isotropic frame and \(k\) for the
positive complementary rank, so the Hermitian dimension is \(2r+k+1\).
Let \(K_0\) be the stabilizer of the completed reference lattice and
let \(K=K_2\), with its given character \(\chi=\chi_2\), be the
original smaller level.  Set \(v_K=\mu_f(K)\).  The level modulus is
the original product \(g=g_0g_1\) of two distinct rational primes.

For a word \(u\), retain its negative vector \(z_u\), its
hyperbolic frame \(e_i,f_i\), its finite transport \(x_u\), its
compact-type space \(V_u\), and its wave function
\(W_u:G_\infty\to V_u\).  The original stabilizer is
\(S_u=N_u\rtimes H_u\).  Here \(H_u\) fixes every
\(e_i,f_i,z_u\), while \(N_u\) acts upper unitriangularly on the
\(e\)-flag and trivially on \(e^\perp/\langle e\rangle\).
The wave-function law, with inner products conjugate linear in their
first variable, is
\[
 W_u(sx)=\psi_{u,\infty}(s)\pi_u(s)W_u(x),
 \qquad \pi_u|_{N_u}=1.
\tag{PGR.1}
\]
The infinite representation is unitary on \(S_u\); \(H_{u,\infty}\)
is a product of compact unitary groups.

The finite set \(I_u\) and representatives \(h_{u,\rho}\) are
those of \(Y.cs(u)\).  They represent
\[
 H_u(F)\backslash H_u(\mathbb A_f)/K_{H,u},\qquad
 K_{H,u}=H_u(\mathbb A_f)\cap x_uKx_u^{-1}.
\]
Put
\[
 a_{u,\rho}=h_{u,\rho}x_u,\qquad
 \widehat\Lambda_{u,\rho}=a_{u,\rho}\widehat L_0,
 \qquad N_{u,\rho}=N_u(F)\cap\operatorname{Aut}(\Lambda_{u,\rho}).
\tag{PGR.2}
\]
The equality with the actual row lattice is a supplied field of
\(Y.cs(u)\).  All the original row conditions, splitting conditions,
finite reduction data and local models of \(Y\) are retained.
They are not newly constructed by this chapter.

The column space \(C_u\) is the original space of linear maps
\(T:V_u\to\{\hbox{functions on }H_{u,\infty}\times H_u(\mathbb A_f)\}\)
with rational left invariance, right \(K_{H,u}\)-character \(\chi\),
and compact-type covariance.  Its scalar seed is exactly
\[
 \Phi_u(T)(x,n h x_u\kappa)
   =\psi_{u,f}(n)\chi(\kappa)\,T(W_u(x))(1,h),
 \quad n\in N_u(\mathbb A_f),\ h\in H_u(\mathbb A_f),\ \kappa\in K;
\tag{PGR.3}
\]
it is zero outside this finite support.  Independence of the
decomposition in (PGR.3) is proved below.  The original Poincar\'e
series, for the arbitrary representatives \(D_u=Y.Dt(u)\), is
\[
 P_uT(x)=\sum_{\delta\in S_u(F)\backslash\Gamma}
                         \Phi_u(T)(\delta x).
\tag{PGR.4}
\]
The quotient notation in this formula means that the representatives
are multiplied on the left of \(x\), as in the source.

Keep the original strict domains \(F_{N,u}=Y.FN(u)\) and set
\[
 c_u=\frac{v_K\,\mu_{N,u}(F_{N,u})\,\mu_{H,u}(H_{u,\infty})}
                 {\dim_{\mathbb C}V_u},\qquad
 s_{u,u'}=\bigl(\sqrt{c_uc_{u'}}\bigr)^{-1}.
\tag{PGR.5}
\]
At the exact formal interface, each mass in this display means its
original real-valued mass, and inverse and division are total, with
\(0^{-1}=0\).  Thus (PGR.5) means the literal
\texttt{c0X}/\texttt{scl}, even before positivity is supplied.
No argument below divides by \(c_u\) to prove an unnormalized identity.
The positive finite number \(v_K\) will be established independently.

For two words, the precise retained analytic assumptions are these.
Both seeds have the original \texttt{FoldOK} data, expanded in
PGR12 below.  Both words satisfy \texttt{InHH}: their scalar series
converge absolutely at every point, are independent of the supplied
coset representatives, lie in the original quotient \(L^2\), and
belong to its specified automorphic range.  The original term set has
\texttt{FundExists}, \texttt{OOIntAll}, \texttt{OOFundIndep}, and
\texttt{AmpWellDef}: its orbital domains exist, its prescribed
operator integrands are integrable, and the integrals and
phase-weighted amplitudes have the stated independence of choices.
Finally the original typed operator datum satisfies \texttt{AbsConvT}.
All quantifiers in these premises are retained, including every column
vector and every permitted domain or representative.  The last premise
includes summability of the Hilbert--Schmidt norms of the original term
operators.  Nothing here presupposes a spectral estimate for their sum.

\subsection{The finite phase, the Levi factor and the columns}

% BEGIN PGR1
\paragraph{PGR1. The supplied additive character annihilates the inverse different.}
\label{claim:PGR1}
Let \(\mathfrak d=\mathfrak d_{F/\mathbb Q}\mathcal O_E\), the
original ideal \(\texttt{ddS}(D)\).  The supplied character
\(\psi_{E,f}\) is unitary and locally constant, and
\[
       \psi_{E,f}(t y)=1
       \quad(y\in\mathfrak d^{-1},\ t\in\widehat{\mathcal O}_E).
\tag{PGR.6}
\]
It therefore annihilates the entire completed fractional ideal
\(\widehat{\mathfrak d^{-1}}\).  This includes the full original
\texttt{PsiDd} statement for every allowed supplied character.
\begin{proof}
The defining identities are multiplicative additivity, triviality on
\(\widehat{\mathcal O}_E\), and
\[
 \psi_{E,f}(b)\exp(2\pi i\operatorname{Tr}_{E/\mathbb Q}b)=1
                         \quad(b\in E).
\tag{PGR.7}
\]
The trace in (PGR.7) is exactly the sum over the chosen CM embeddings
and their conjugates in the original infinite character.  For an
arbitrary finite adele \(a\), some positive integer \(m\) makes
\(ma\) integral.  Additivity gives \(\psi_{E,f}(a)^m=1\).
Thus every value has modulus one.  Triviality on the open integral
subring makes the character locally constant; no continuity assumption
has been inserted.

If \(y\in\mathfrak d^{-1}\) and \(o\in\mathcal O_E\), then
\(\operatorname{Tr}_{E/F}(oy)\in\mathfrak d_{F/\mathbb Q}^{-1}\).
Indeed multiplication by each element of \(\mathfrak d_{F/\mathbb Q}\)
makes \(oy\) integral over \(\mathcal O_F\), and hence makes its
relative trace integral.  The defining trace-dual property of the
different now gives
\(\operatorname{Tr}_{E/\mathbb Q}(oy)\in\mathbb Z\).

Choose a nonzero integer \(m\) with \(my\in\mathcal O_E\).
Finite Chinese remaindering gives
\(t=o+mt_1\), with \(o\in\mathcal O_E\) and
\(t_1\in\widehat{\mathcal O}_E\).  Consequently
\(ty=oy+t_1(my)\).  The second term is integral, and the first
has integral absolute trace.  Equation (PGR.7) proves (PGR.6).
A finite set of integral-module generators for
\(\mathfrak d^{-1}\) generates its completed module as well;
apply (PGR.6) to each generator and multiply the character values.
This proves the last assertion, including at ramified finite places.
\end{proof}
% END PGR1

% BEGIN PGR2
\paragraph{PGR2. Rational splitting and the actual unipotent characters.}
\label{claim:PGR2}
Over \(F\), every completion of \(F\), and the finite adeles,
the actual group \(S_u\) is \(N_u\rtimes H_u\).  The splitting
is defined over \(F\).  The characters \(\psi_{u,f}\) and
\(\psi_{u,\infty}\) are unitary characters of \(S_u\), trivial
on \(H_u\), invariant under its conjugation on \(N_u\), and
\[
       \psi_{u,f}(n_f)\psi_{u,\infty}(n_\infty)=1
                           \quad(n\in N_u(F)).
\tag{PGR.8}
\]
In particular this proves the original rational Levi and rational
character-product assertions.
\begin{proof}
Let \(U=\langle e_0,\ldots,e_{r-1}\rangle\), and let \(W\)
be the complement orthogonal to all \(e_i,f_i\).  The hyperbolic
pairings give the rational projection
\[
 p_W(y)=y-\sum_i h(f_i,y)e_i-\sum_i h(e_i,y)f_i.
\tag{PGR.9}
\]
An element of \(S_u\) induces an isometry of \(U^\perp/U\),
identified with \(W\), fixing \(z_u\).  Extend this isometry
by the identity on every \(e_i,f_i\); the result is its unique
Levi factor in \(H_u\).  The residual factor is exactly in
\(N_u\).  This construction and its inverse use the displayed
projection and the unitary equations, so commute with all the stated
scalar extensions.  In particular a rational element has rational
factors.

In the original character put
\[
 B_u(s)=\sum_{i+1<r}h(f_i,se_{i+1})+
             \begin{cases}h(f_{r-1},sz_u),&r>0,\\0,&r=0.\end{cases}
\tag{PGR.10}
\]
The finite character is \(\psi_{E,f}(B_u(s)/g)\); the infinite
one is the product of the corresponding real exponential characters.
The upper unitriangular action on the \(e\)-flag gives
\(B_u(st)=B_u(s)+B_u(t)\): in an adjacent entry there is no
intermediate index, and in the last row of the \(z_u\)-column
all earlier-row contributions vanish.  A Levi factor fixes all the
vectors occurring in (PGR.10).  Unitarity then gives
\(B_u(hnh^{-1})=B_u(n)\) and \(B_u(h)=0\).
This proves the character and conjugacy statements.  Their values
are unitary by PGR1 and the explicit real exponentials.
For rational \(n\), (PGR.10) is a rational element of \(E\),
and (PGR.8) is precisely (PGR.7) applied to \(B_u(n)/g\).
All sums are empty, with value zero, when \(r=0\).
\end{proof}
% END PGR2

% BEGIN PGR3
\paragraph{PGR3. The exact finite-level phase, including the different.}
\label{claim:PGR3}
For every original \(X,u\) and every
\(n\in N_u(\mathbb A_f)\cap K_0\),
\[
 x_u^{-1}nx_u\in K,\qquad
             \chi(x_u^{-1}nx_u)=\psi_{u,f}(n).
\tag{PGR.11}
\]
This is the complete original \texttt{UnipotentConjugateInLevel}.
\begin{proof}
The membership assertion is the already proved RCV12, whose scope is
every supplied finite reference frame and transport.  It imposes no
character comparison.  We prove that remaining comparison here.
Let \(A\) be the actual numerator defining
\(\chi(x_u^{-1}nx_u)=\psi_{E,f}(A/g)\), including both its
\(e,z\) chain and its \(bw\) chain.  Let \(B=B_u(n)\).
All vectors in the numerator \(A\) are integral, and
\(x_u^{-1}nx_u\) preserves the integral self-dual lattice.
Thus \(A\in\widehat{\mathcal O}_E\).

We first locate \(B\) at every finite place.  The actual lattice
splits as
\[
 L_0=\bigoplus_{i<r}\mathfrak a_i e_i\ \oplus L_b\ \oplus
               \bigoplus_{i<r}\mathfrak a_i^{-1} f_i,
             \qquad \mathfrak a_i=\mathfrak d^{i-r}.
\tag{PGR.12}
\]
Projection onto its \(e_i\)-coordinate sends it into
\(\mathfrak a_i\).  Since \(n\) preserves the completed
lattice, its coefficient of \(e_i\) in \(ne_{i+1}\) belongs
to \(\mathfrak a_i\mathfrak a_{i+1}^{-1}=\mathfrak d^{-1}\).
Also \(z_u\in L_b\), so the coefficient of \(e_{r-1}\)
in \(nz_u\) belongs to \(\mathfrak a_{r-1}=\mathfrak d^{-1}\).
These are exactly the summands of \(B\).  Therefore
\(B\in\widehat{\mathfrak d^{-1}}\).

At each place above \(g\), the supplied reference vectors agree
modulo \(g\) with the rational vectors, after the actual transport:
\(x_ue_i^0\equiv e_i\), \(x_uf_i^0\equiv f_i\), and
\(x_uz^0\equiv z_u\).  These congruences have precisely the
localized meaning proved in ASG11; their clearing integers are units
at these places.  Applying the lattice isometries \(n,x_u\) and
pairing preserves the congruences.  The \(e,z\) part of \(A\)
is thus congruent to \(B\) modulo \(g\).  The strengthened
\(bw\)-congruence in the proof of RCV12 says that
\((x_u^{-1}nx_u-1)bw_j\) is in the span of the \(e_i^0\)
modulo \(g\).  Pairing with \(bwd_i\) annihilates this span
modulo \(g\), and the adjacent pairing
\(h(bwd_i,bw_{i+1})\) is zero modulo \(g\).  Hence the
entire additional \(bw\) numerator is zero modulo \(g\).
We obtain \((A-B)/g\in\widehat{\mathcal O}_{E,v}\) there.

Away from \(g\), division by \(g\) is an integral unit.
Since \(\mathcal O_E\subset\mathfrak d^{-1}\), the preceding
locations of \(A\) and \(B\) show that \((A-B)/g\) belongs
to the local completed inverse different there too.  Thus globally
\((A-B)/g\in\widehat{\mathfrak d^{-1}}\).  PGR1 gives
\(\psi_{E,f}((A-B)/g)=1\), proving (PGR.11).  This argument
explains why integrality of the rational dual vectors at every place
is neither assumed nor needed.
\end{proof}
% END PGR3

% BEGIN PGR4
\paragraph{PGR4. Integral Levi extraction and decomposition independence.}
\label{claim:PGR4}
If \(n\in N_u(\mathbb A_f)\), \(h\in H_u(\mathbb A_f)\),
and \(nh\in K_0\), then both \(h,n\in K_0\).
Moreover
\[
 S_u(\mathbb A_f)\cap x_uKx_u^{-1}
   =(N_u(\mathbb A_f)\cap K_0)\rtimes K_{H,u}.
\tag{PGR.13}
\]
The displayed factorization is unique.  These assertions supply the
original \texttt{KsetSplit} and its smaller-level consequence.
\begin{proof}
The projection (PGR.9) preserves the completed lattice by (PGR.12):
it is zero on its two isotropic summands and the identity on \(L_b\).
For an integral vector \(y\), put \(p=p_W(y)\).  The Levi
factor preserves \(W\), and \(n\) is the identity on
\(U^\perp/U\), so
\[
                  p_W(nhp)=hp.
\]
The left side is integral because \(p\) is integral and \(nh\)
preserves the lattice.  The vector \(y-p\) lies in the hyperbolic
summand and is fixed by \(h\).  Hence
\(hy=(y-p)+hp\) is integral.  Apply the same argument to
\((nh)^{-1}=(h^{-1}n^{-1}h)h^{-1}\); PGR2 ensures that
its first factor lies in \(N_u\).  It follows that \(h^{-1}\)
also preserves the lattice.  Thus \(h\in K_0\), and then
\(n=(nh)h^{-1}\in K_0\).

If \(nh\in x_uKx_u^{-1}\), it is in \(K_0\) because
\(x_u\in K_0\).  PGR3 puts \(n\) in \(x_uKx_u^{-1}\).
Dividing gives \(h\in K_{H,u}\).  Conversely both factors on
the right of (PGR.13) belong to the left by PGR3.  Uniqueness is
the uniqueness of the splitting in PGR2.  Notice that the proof uses
the actual fractional-ideal splitting, not a globally integral
hyperbolic basis that the data do not provide.
\end{proof}
% END PGR4

% BEGIN PGR5
\paragraph{PGR5. Triviality of every original column stabilizer.}
\label{claim:PGR5}
For every \(h\in H_u(\mathbb A_f)\),
\[
          H_u(F)\cap hK_{H,u}h^{-1}=\{1\}.
\tag{PGR.14}
\]
This is the full original \texttt{NeatKu}, for every word and
every permitted level-character witness.
\begin{proof}
Let \(\gamma\) belong to the intersection.  It preserves the
rational lattice whose completion is \(hx_u\widehat L_0\).
This is a full lattice: clearing the finitely many finite-place
denominators bounds it between \(mL_0\) and \(m^{-1}L_0\)
for some positive integer \(m\), and intersection with the rational
space recovers the stated completed lattice.  The element \(\gamma\)
fixes the hyperbolic pairs and \(z_u\).  Their further orthogonal
complement is positive definite at every real place.  On that
complement, the absolute trace of the Hermitian form is a positive
definite rational quadratic form.  Project the full lattice onto
this rational complement.  The projection is rational, so its image
is a full integral lattice after clearing one denominator.
The induced automorphism preserves that lattice and its positive
definite quadratic form.

Such a lattice has only finitely many isometries: the image of each
fixed integral basis vector lies in a bounded set of lattice points,
and there are finitely many possibilities for all basis images.
The restriction is faithful because the other summands are fixed.
Consequently \(\gamma\) has finite order.

The integral matrix
\(\kappa=x_u^{-1}h^{-1}\gamma_fhx_u\) belongs to \(K\).
The original residue-frame calculation in ASG14 makes its reduction
unipotent at both \(g_0\) and \(g_1\).  To recall why no rank
restriction is hidden here, the \(e\)-flag is upper unitriangular;
the \(z\)-column is fixed modulo that flag; and the \(bw\)-flag
is upper unitriangular after projection off those fixed summands.
The integral idempotent and its conjugate split the remaining
residue module.  Their dual pairings make the action on the conjugate
summand the inverse dual.  It too is unipotent.  The supplied spanning
condition accounts for the whole residue module.

If an integral finite-order matrix has unipotent reduction in residue
characteristic \(p\), its order is a power of \(p\).  Indeed write
its order as \(p^ab\), with \((p,b)=1\).  Its \(p^a\)-th
power \(M\) has order dividing \(b\); its reduction is unipotent
and has order dividing \(b\), hence is the identity.  In
\((M-I)(I+M+\cdots+M^{b-1})=0\), the second factor reduces
to the invertible matrix \(bI\), and is invertible over the local
ring.  Thus \(M=I\).  Applying this at the two distinct primes
forces the order of \(\gamma\) to be one.  This works at ramified
places as well and never compares a prime with the rank.
\end{proof}
% END PGR5

% BEGIN PGR6
\paragraph{PGR6. The original column coordinates are antilinear Hilbert coordinates.}
\label{claim:PGR6}
For \(T\in C_u\), let \(v_{u,\rho}(T)\in V_u\) be determined by
\[
       T(\xi)(1,h_{u,\rho})
                    =\langle v_{u,\rho}(T),\xi\rangle.
\tag{PGR.15}
\]
Then \(J_uT=(v_{u,\rho}(T))_\rho\) is exactly the original
\texttt{sharp}; it is an antilinear bijection onto
\(\mathcal V_u=\bigoplus_{\rho\in I_u}V_u\).  The original
\texttt{ColT}/\texttt{colCore} inner product is
\[
       \langle T,T'\rangle_{C_u}=\langle J_uT',J_uT\rangle_{\mathcal V_u}.
\tag{PGR.16}
\]
Thus \(J_u\) is an antilinear isometry; in particular the original
column space is finite dimensional with this prescribed norm.
\begin{proof}
Riesz representation on the finite-dimensional space \(V_u\)
gives (PGR.15), additivity and conjugate scalar multiplication.
Every \(h_f\) has a unique column class and a presentation
\(h_f=\gamma_fh_{u,\rho}\kappa\), with
\(\gamma\in H_u(F)\), \(\kappa\in K_{H,u}\).
For a fixed class, PGR5 makes \(\gamma\), and therefore
\(\kappa\), unique.  An arbitrary list of linear functionals
\(\lambda_\rho\in V_u^*\) consequently defines
\[
 T(\xi)(h_\infty,h_f)
   =\chi(x_u^{-1}\kappa x_u)\,
              \lambda_\rho\bigl(\pi_u(\gamma_\infty^{-1}h_\infty)\xi\bigr).
\tag{PGR.17}
\]
Uniqueness of the presentation proves this is well defined.  The
representation and character laws verify rational left invariance,
finite right covariance and compact-type covariance directly.
Conversely those three laws force (PGR.17) from the evaluations
(PGR.15).  Hence evaluation is a linear bijection with the list of
duals, and Riesz representation makes \(J_u\) an antilinear
bijection.  Equation (PGR.16) is the literal transported inner product,
not a new assertion about an earlier norm on a function space.
It is positive definite by injectivity, and its norm is
\(\|T\|^2=\sum_\rho\|v_{u,\rho}(T)\|^2\), as claimed.
\end{proof}
% END PGR6

% BEGIN PGR7
\paragraph{PGR7. The original scalar seed is well defined and has the required covariance.}
\label{claim:PGR7}
Formula (PGR.3) is independent of its decomposition.  It defines a
measurable function, invariant under \(S_u(F)\) on the left and
transforming by \(\chi\) under \(K\) on the right.  It also
satisfies the original convolution reproduction formula.  Thus both
\texttt{PhWellDef} and \texttt{PhProps} hold on every supplied \(X\).
\begin{proof}
Suppose \(nhx_u\kappa=n'h'x_u\kappa'\).  The element
\((n'h')^{-1}nh\) belongs to \(x_uKx_u^{-1}\).
Its Levi factor is \(h'^{-1}h\), and PGR4 says this lies in
\(K_{H,u}\); its unipotent factor lies in \(N_u\cap K_0\).
The character equality in PGR3, the conjugation invariance in PGR2,
and the right covariance of \(T\) then give equality of the two
values in (PGR.3).  More explicitly, write
\(n'^{-1}n=h' n_0 h'^{-1}\) and \(h=h'h_0\).  Then
\(n_0h_0=x_u\kappa'\kappa^{-1}x_u^{-1}\), and
\[
 \psi_{u,f}(n)=\psi_{u,f}(n')\psi_{u,f}(n_0),\qquad
 T(\xi)(1,h)=\chi(x_u^{-1}h_0x_u)T(\xi)(1,h').
\]
Their product with \(\chi(\kappa)\) is the expression using
\(n',h',\kappa'\), by (PGR.11).

Right \(K\)-covariance is now immediate.  For left rational
covariance split \(s=n_sh_s\in S_u(F)\) by PGR2.
At infinity (PGR.1) contributes
\(\psi_{u,\infty}(n_s)\pi_u(h_s)\).  At finite places the
unipotent factor contributes \(\psi_{u,f}(n_s)\), while
rational left invariance and compact-type covariance of \(T\)
remove the common \(h_s\).  Equation (PGR.8) cancels the two
remaining factors.  The support is stable under these operations,
so the same assertions hold at points where the seed is zero.

The finite support is a union of right \(K\)-cosets.  Each coset
has a chosen decomposition, and on it (PGR.3) is a continuous
function of \(x\) times a locally constant finite character.
The group \(G_f/K\) is countable: a second-countable space has
at most countably many disjoint open cosets.  Thus the piecewise
formula is measurable; its value outside the union is zero.
Here \(K\) is compact open, with positive finite mass, by the
universal residue-frame argument ASG14.  That argument applies
because PGR1 supplies exactly the original character properties;
it is independent of a choice of construction for \(X\).

For a fixed finite component, the seed is a scalar multiple of a
continuous linear functional applied to \(W_u(x)\).
Apply that functional to the Bochner-integrable reproduction formula
in \(X.Wu\_spec\).  A continuous linear functional commutes with
the integral and preserves integrability.  This gives precisely the
scalar reproduction clause of \texttt{PhProps}, including zero
outside the finite support.  No assertion that the reproducing
kernel is compactly supported is needed here.
\end{proof}
% END PGR7

% BEGIN PGR8
\paragraph{PGR8. Rational unipotent approximation and strict column domains.}
\label{claim:PGR8}
The subgroup \(N_u(F)\) is dense in \(N_u(\mathbb A_f)\).
Consequently \(N_u(\mathbb A_f)=N_u(F)U_f\) for every open
subgroup \(U_f\) of \(N_u(\mathbb A_f)\).  Each
\(N_{u,\rho}\), in its original infinite embedding, is countable
and discrete and admits a measurable strict fundamental domain
\(\Omega_{u,\rho}\) in \(G_\infty\).
\begin{proof}
In a basis adapted to the hyperbolic flag, the unipotent Lie algebra
is a finite-dimensional \(F\)-vector space of strictly triangular
matrices satisfying linear skew-Hermitian equations.  Exponential
and logarithm are finite polynomials, inverse to each other on this
space and \(N_u\).  To check this, logarithm respects the flag and
is zero on the successive diagonal quotients.  The unitary relation
\(g^*=g^{-1}\) gives \((\log g)^*=-\log g\); conversely the
exponential of a skew-Hermitian nilpotent matrix is unitary and has
the required flag and quotient action.  The finite polynomials are
defined over \(F\) and are continuous at every place.
Equivalently, the complete polynomial reconstruction
in RCV14 gives these same rational and adelic points.  Finite adelic
approximation in a rational basis of the Lie algebra therefore proves
density.  An open coset \(nU_f\) meets \(N_u(F)\), proving the
factorization assertion.

The group \(N_{u,\rho}\) preserves the full row lattice.
Its matrix entries in an integral basis lie in a fixed discrete
integer lattice inside the product of the real matrix spaces.
They are therefore discrete at infinity, and countable.  More
generally any discrete subgroup \(L\) of a second-countable locally
compact group has a measurable strict domain for left translation.
To see this, choose a countable cover by open sets so small that
distinct \(L\)-translates of each are disjoint.  From the first
open set take all its points; from each later set remove all translates
of the points already taken.  The remaining pieces are Borel, meet
each orbit once, and their countable union is the required domain.
Apply this construction in \(G_\infty\).

The same elementary reasoning supplies fundamental domains for the
countable rational subgroups needed in adelic unfolding.  The
diagonal rational group is discrete in \(G_{\mathbb A}\): bound
finite denominators in a compact finite neighborhood, then use the
discreteness of its integral coordinates in all infinite embeddings.
The supplied quotient domain \(X.B.fund\) already has the original
almost-everywhere fundamental-domain property, which is sufficient
for every integral below.  We neither replace it nor normalize its
measure anew.
\end{proof}
% END PGR8

% BEGIN PGR9
\paragraph{PGR9. The exact column unipotents and their phases.}
\label{claim:PGR9}
For every rational \(n\in N_u(F)\) and every column \(\rho\),
\[
 n\in N_{u,\rho}\quad\Longleftrightarrow\quad
             a_{u,\rho}^{-1}n_fa_{u,\rho}\in K.
\tag{PGR.18}
\]
When these conditions hold,
\[
 \chi(a_{u,\rho}^{-1}n_fa_{u,\rho})=\psi_{u,f}(n_f)
                              =\psi_{u,\infty}(n_\infty)^{-1}.
\tag{PGR.19}
\]
Thus the subgroup here is exactly the source's \texttt{NOcol},
with its unique rational lift, not a finite-index replacement.
\begin{proof}
Preservation of \(\Lambda_{u,\rho}\) by a rational isometry is
equivalent to preservation of its completion.  This follows by
expressing a rational vector in an integer basis: being integral
in every finite completion is equivalent to having integer rational
coordinates.  The converse for the completion follows by scalar
extension.  Since \(x_u\) preserves \(\widehat L_0\), the
left side of (PGR.18) says
\(h_{u,\rho}^{-1}n_fh_{u,\rho}\in K_0\).
This is an element of \(N_u\) by PGR2.  PGR3 proves the required
membership in \(K\) and the first equality in (PGR.19), using
invariance of the unipotent character under \(h_{u,\rho}\).
Conversely \(K\subset K_0\) proves lattice preservation.
The last equality is (PGR.8).  The rational embedding into either
the finite adeles or all infinite places is injective, so there is
exactly one rational lift of the infinite element.  The original
definition of \(N\!\operatorname{Lam}\) is unipotent flag
preservation together with lattice preservation, giving the claimed
literal identification.
\end{proof}
% END PGR9

% BEGIN PGR10
\paragraph{PGR10. The support cells have exactly the original product measure.}
\label{claim:PGR10}
For arbitrary strict domains \(\Omega_{u,\rho}\) as in PGR8,
the sets
\[
              \Omega_{u,\rho}\times a_{u,\rho}K
\tag{PGR.20}
\]
form a disjoint set of representatives for the support of the seed in
\(S_u(F)\backslash G_{\mathbb A}\).  If a measurable
\(S_u(F)\)-invariant function \(f\) is supported there, its
quotient integral is the sum of its integrals over (PGR.20), first
for nonnegative \(f\), and then for every integrable complex \(f\).
If \(f\) is also right \(K\)-invariant, each summand is
\[
                  v_K\int_{\Omega_{u,\rho}}f(x,a_{u,\rho})\,d\mu_\infty(x).
\tag{PGR.21}
\]
\begin{proof}
For coverage, start with a finite component \(nhx_u\kappa\).
Reduce \(h\) by its supplied rational Levi double class, then
use PGR8 with
\(U_f=N_u(\mathbb A_f)\cap a_{u,\rho}Ka_{u,\rho}^{-1}\)
to remove its finite unipotent component by a rational unipotent.
The finite component is now in \(a_{u,\rho}K\).  Finally move
its infinite component into \(\Omega_{u,\rho}\) using
\(N_{u,\rho}\); (PGR.18) keeps the finite component in that
same coset.

For uniqueness, suppose a rational \(s\in S_u(F)\) carries
a point of the \(\rho\)-cell to a point of the \(\sigma\)-cell.
Write \(s=nh\) by PGR2.  The equality of finite components,
conjugated by the two \(h\)-representatives, puts the resulting
unipotent--Levi product in \(x_uKx_u^{-1}\).
PGR4 shows that its Levi factor lies in \(K_{H,u}\).
Uniqueness of the original column classes forces \(\rho=\sigma\).
The remaining rational Levi element stabilizes that column, hence
is the identity by PGR5.  The residual \(s\) is in
\(N_{u,\rho}\), by (PGR.18); strictness of the infinite domain
then gives \(s=1\).

To prove the integral statement, intersect (PGR.20) with the
countably many translates of any supplied adelic fundamental domain.
The pieces are null measurable and partition both systems of
representatives, up to the allowed null sets.  For integration replace
the supplied adelic domain by a measurable representative; its difference
and the countably many Haar translates of that difference are null.
This is the measurable-representative convention detailed in PGR11.
Left Haar invariance and invariance of
\(f\) identify their integrals.  Tonelli proves the assertion for
nonnegative functions without a prior finiteness assumption; apply
it to \(|f|\) before making the complex-valued change of domain.
Finally \(\mu_{\mathbb A}=\mu_\infty\times\mu_f\), and
\(\mu_f(a_{u,\rho}K)=v_K\).  Product integration and finite
right invariance give (PGR.21).  A rational unipotent also translates
the point of \(K\); this causes no extra factor because \(\mu_f\)
is left invariant.  There is no column-stabilizer denominator, precisely
because PGR5 proved it is one.
\end{proof}
% END PGR10

\subsection{Absolute unfolding before the orbital series is summed}

% BEGIN PGR11
\paragraph{PGR11. The full folding inequality and its compact finiteness bound.}
\label{claim:PGR11}
Let \(G\) be a second countable locally compact topological group,
\(\mu\) a left Haar measure, \(\Gamma\) a discrete subgroup,
and \(S\leq\Gamma\).  Let \(D\) be any left-coset
representatives for \(S\backslash\Gamma\), and \(F_S\) any
original almost-everywhere fundamental domain for the left action
of \(S\).  Thus \(F_S\) need only be null measurable; covering
and disjointness of its translates are allowed to hold almost everywhere.
For a measurable set \(\Omega\), put
\[
 m_\Omega(x)=\sup_{y\in G}
       \#\{\gamma\in\Gamma:x^{-1}\gamma y\in\Omega\}.
\tag{PGR.22}
\]
Counts and the supremum may be infinite.  If a measurable,
left \(S\)-invariant scalar function \(\phi\) satisfies
\[
 |\phi(y)|\leq a\int_\Omega|\phi(y\omega)|\,d\mu(\omega),
 \qquad a\geq0,
\tag{PGR.23}
\]
then, as an inequality of extended nonnegative integrals,
\[
 \sum_{\delta\in D}|\phi(\delta x)|
       \leq a\,m_\Omega(x)\int_{F_S}|\phi(y)|\,d\mu(y).
\tag{PGR.24}
\]
This is the full original \texttt{FoldBound}, including arbitrary
measurable \(\Omega\), without a Hausdorff assumption.
If in addition \(G\) is Hausdorff and \(\Omega\) is compact
and contains \(1\), then \(m_\Omega\) is finite, Borel measurable and at
least one.  In particular, with
\[
 A_\phi=\int_{F_S}|\phi|,\qquad
 J_\phi=\int_{F_S}m_\Omega|\phi|<\infty,
\tag{PGR.25}
\]
we have \(A_\phi\leq J_\phi\) and the sum in (PGR.24)
is finite at every \(x\).
\begin{proof}
A discrete subgroup of a second countable group is countable.
The domain convention causes no extra regularity assumption.
Replace each null-measurable domain by a measurable representative
modulo \(\mu\).  Every Haar translate of the difference is null.
The union of these differences, the exceptional set for coverage,
and the intersections witnessing failure of disjointness is still
null, because the acting groups are countable.  The translated
representatives therefore give the required partitions off one null
set, with all original set integrals unchanged.  The domain partitions
below, and in PGR12, are understood in this sense.  This does not
weaken the separate strict \texttt{IsSFund} statements in this chapter.
Apply (PGR.23) at \(\delta x\), and change variable
\(y=\delta x\omega\), which uses only left invariance of
\(\mu\).  Sum the resulting nonnegative integrals.  Partition
\(G\) by \(sF_S\), \(s\in S\), and use the invariance
of \(|\phi|\).  The result is at most
\[
 a\int_{F_S}\sum_{\delta\in D,\ s\in S}
    1_\Omega(x^{-1}\delta^{-1}s y)\,|\phi(y)|\,d\mu(y).
\tag{PGR.26}
\]
Every \(\gamma\in\Gamma\) has exactly one expression
\(\delta^{-1}s\): apply the given coset decomposition to
\(\gamma^{-1}\).  Thus the sum of indicators in (PGR.26)
is the count occurring in (PGR.22), for this particular \(y\),
and is bounded by \(m_\Omega(x)\).  This proves (PGR.24).
All interchanges so far are Tonelli interchanges of nonnegative
measurable functions.  No integrability or right invariance of Haar
measure was required.

For the remaining assertions assume \(G\) is Hausdorff and
\(\Omega\) compact.  The discrete subgroup \(\Gamma\) is
then closed and meets compact sets in finite sets.
If one element \(\gamma_0\) is
counted at \((x,y)\), every other counted element satisfies
\[
 \gamma\gamma_0^{-1}\in
                     \Gamma\cap x\Omega\Omega^{-1}x^{-1}.
\]
The latter set is finite and is independent of \(y\).  Its
cardinality bounds \(m_\Omega(x)\).  Choosing \(y=x\)
and \(\gamma=1\) proves \(m_\Omega(x)\geq1\).
For completeness, its measurability can be seen without taking an
uncontrolled supremum of Borel functions.  The condition
\(m_\Omega(x)\geq q\), for a positive integer \(q\), is the
union over the countably many distinct tuples
\(\gamma_1,\ldots,\gamma_q\) of the conditions that some
\(\omega\in\Omega\) satisfies
\[
       x^{-1}\gamma_i\gamma_1^{-1}x\omega\in\Omega
                                  \quad(1\leq i\leq q).
\]
For each tuple this is closed: it is the projection of a closed subset
of \(G\times\Omega\), with compact second factor.  The union
is Borel.  This proves the stated measurability.  The final assertions
follow from (PGR.24), (PGR.25), and \(m_\Omega\geq1\).
\end{proof}
% END PGR11

% BEGIN PGR12
\paragraph{PGR12. The original mixed Poincar\'e identity, with absolute control.}
\label{claim:PGR12}
In the compact Hausdorff situation of PGR11, let \(\phi'\) be a
left \(S'\)-invariant seed satisfying the same conditions with
constant \(a'\), domain \(F_{S'}\), and finite quantities
\(A_{\phi'},J_{\phi'}\).  Use any respective representatives
\(D,D'\), and write
\[
 P\phi(x)=\sum_{\delta\in D}\phi(\delta x),\qquad
 P^{\rm abs}\phi(x)=\sum_{\delta\in D}|\phi(\delta x)|.
\]
Both series converge at every point and are independent of the
representatives.  For every \(\Gamma\)-fundamental domain
\(F_\Gamma\),
\[
 \begin{split}
 \int_{F_S}|\phi|P^{\rm abs}\phi'
       &\leq a'A_{\phi'}J_\phi<\infty,\\
 \int_{F_{S'}}|\phi'|P^{\rm abs}\phi
       &\leq aA_\phi J_{\phi'}<\infty,\\
 \int_{F_\Gamma}|P\phi|^2
       &\leq a A_\phi J_\phi\leq aJ_\phi^2.
 \end{split}
\tag{PGR.27}
\]
The functions \(P\phi,P\phi'\) are measurable and belong to
the indicated quotient \(L^2\).  If \(R\) is any double-coset
representative set for \(S\backslash\Gamma/S'\), put
\(S_r=S\cap rS'r^{-1}\), and choose arbitrary fundamental
domains \(F_r\) for these groups, with the same null-measurable,
almost-everywhere convention as in PGR11.  Then
\[
 \sum_{r\in R}\int_{F_r}|\phi(x)|\,|\phi'(r^{-1}x)|\,d\mu(x)
                                                   <\infty,
\tag{PGR.28}
\]
and
\[
 \begin{split}
 \int_{F_\Gamma}\overline{P\phi(x)}P\phi'(x)\,d\mu(x)
   &=\int_{F_S}\overline{\phi(x)}P\phi'(x)\,d\mu(x)\\
   &=\sum_{r\in R}\int_{F_r}
          \overline{\phi(x)}\phi'(r^{-1}x)\,d\mu(x).
 \end{split}
\tag{PGR.29}
\]
Taking complex conjugates gives the source's other inner-product
orientation.  Thus this proves the full original
\texttt{PoincareSeriesFinite}, \texttt{PoincareSeriesMemLp}, and
\texttt{PoincareSeriesInner} at their stated general group scopes.

For the actual setting, \texttt{FoldOK} supplies a domain \(F_S\)
and a compact \(\Omega_0\ni1\), with one constant \(a\)
before all column vectors \(T\).  It supplies (PGR.23) for every
compact \(\Omega\supset\Omega_0\), and the finiteness of
\(J_{\Phi_u(T)}\) for every compact \(\Omega\) and every
\(S_u(F)\)-fundamental domain.  For two words we take the one
common compact set \(\Omega_0\cup\Omega'_0\).  All the
conclusions above therefore apply to their seeds for every \(T,T'\).
\begin{proof}
PGR11 gives pointwise absolute convergence.  Changing a representative
multiplies it on the left by an element of the seed's stabilizer,
which changes no summand.  Right multiplication by \(\Gamma\)
permutes the left cosets.  Consequently both periodized functions,
including their absolute versions, are left \(\Gamma\)-invariant.
They are measurable as countable sums of measurable functions.

Multiply the pointwise bound
\(P^{\rm abs}\phi'\leq a'm_\Omega A_{\phi'}\)
by \(|\phi|\) and integrate over \(F_S\).  This is the
first inequality of (PGR.27); interchange the seeds for the second.
To unfold, partition each translate of \(F_\Gamma\) by its
intersection with \(F_S\), or equivalently tile an
\(S\)-domain by the translates \(\delta F_\Gamma\),
\(\delta\in D\).  For every nonnegative left
\(\Gamma\)-invariant measurable \(f\), this gives
\[
 \int_{F_\Gamma}P^{\rm abs}\phi(x) f(x)\,d\mu(x)
                 =\int_{F_S}|\phi(x)|f(x)\,d\mu(x).
\tag{PGR.30}
\]
This is a nonnegative identity before it is an integrability assertion.
Apply it with \(f=P^{\rm abs}\phi'\).  The first bound of
(PGR.27) proves that all scalar products to be unfolded are dominated
by an integrable function.  Applying the same reasoning to a seed
paired with itself proves the third bound, hence its \(L^2\)
membership.  Dominated integration of the now absolutely convergent
scalar sums proves the first identity in (PGR.29).

For the second identity, replace the left-coset index \(\delta'\)
by its inverse, a right-coset index in \(\Gamma/S'\).  The left
action of \(S\) on these right cosets has orbits indexed by
\(R\), and the stabilizer of \(rS'\) is \(S_r\).
Choose representatives \(s\) of \(S/S_r\).  The corresponding
summand is \(\phi'((sr)^{-1}x)\).  In its integral over
\(F_S\), set \(y=s^{-1}x\); left \(S\)-invariance of
\(\phi\) gives the integrand
\(\overline{\phi(y)}\phi'(r^{-1}y)\).
The sets \(s^{-1}F_S\), for these representatives, form an
\(S_r\)-domain.  Partitioning intersections with the arbitrary
given \(F_r\) changes neither the nonnegative integral nor the
integrable scalar integral.  First doing this with absolute values
identifies the sum in (PGR.28) with the first integral in (PGR.27).
It is finite.  Every term is therefore integrable, and the same
partition proves the complex identity.  This also explains why an
arbitrary choice of the domains \(F_r\) is permitted.

Finally, the definition of \texttt{FoldOK} has precisely the
uniformity stated above.  The union of the two supplied compact sets
is compact, contains one, and permits both averaging inequalities and
both weighted integrability bounds.  No pointwise decay exponent or
arithmetic restriction on the words has been used.
\end{proof}
% END PGR12

% BEGIN PGR13
\paragraph{PGR13. Absolute slab integrability and single unfolding in the actual measure.}
\label{claim:PGR13}
For every pair of column vectors \(T,T'\) and every family of
strict domains \(\Omega_{u,\rho}\) for \(N_{u,\rho}\),
\[
 \int_{\Omega_{u,\rho}}|\Phi_u(T)(x,a_{u,\rho})|
       \sum_{\delta\in D_{u'}}
          |\Phi_{u'}(T')(\delta(x,a_{u,\rho}))|
                     \,d\mu_\infty(x)<\infty.
\tag{PGR.31}
\]
Thus the full original \texttt{SlabAbs} holds on the supplied
setting.  In the actual Hilbert space selected by \texttt{InHH},
\[
 \langle P_uT,P_{u'}T'\rangle
   =v_K\sum_{\rho\in I_u}\int_{\Omega_{u,\rho}}
       \overline{\langle J_uT(\rho),W_u(x)\rangle}
          P_{u'}T'(x,a_{u,\rho})\,d\mu_\infty(x).
\tag{PGR.32}
\]
There is exactly one finite-level factor \(v_K\).
\begin{proof}
Take \(\phi=\Phi_u(T)\), \(\phi'=\Phi_{u'}(T')\)
in PGR12.  The nonnegative function
\(|\phi|P^{\rm abs}\phi'\) is left \(S_u(F)\)-invariant,
measurable, and supported on the first seed's finite support.  Both
factors acquire unit-modulus characters under right multiplication by
\(K\), so their absolute values are right \(K\)-invariant.
Its integral over an \(S_u(F)\)-domain is finite by (PGR.27).
PGR10 identifies that integral with \(v_K\) times the sum of
the nonnegative integrals in (PGR.31).  There are finitely many columns,
and \(0<v_K<\infty\) by ASG14.  Each integral is finite.

Now apply the same change of domain to
\(\overline{\phi}P\phi'\).  Its integrability has just been
proved; its right \(K\)-characters cancel.  PGR6 identifies the
first seed on the cell with
\(\langle J_uT(\rho),W_u(x)\rangle\).
The first identity of (PGR.29) gives (PGR.32).  The quotient measure
here is the original measure obtained from \(X.B.fund\), which is
an allowed \(\Gamma\)-fundamental domain.  The supplied
\texttt{InHH} identifies its measurable scalar periods with the
classes used to define \(P_u\) and \(P_{u'}\) in that Hilbert
space.  Equality is consequently equality of the original inner
products, rather than of an independently normalized quotient model.
\end{proof}
% END PGR13

\subsection{Rational returns and the literal orbital terms}

% BEGIN PGR14
\paragraph{PGR14. The return dictionary includes the full original flag predicate.}
\label{claim:PGR14}
Fix columns \(\rho\in I_u\), \(\rho'\in I_{u'}\), and put
\[
 \Gamma_{\rho,\rho'}=
 \{r\in\Gamma:k_r=a_{u,\rho}^{-1}r_fa_{u',\rho'}\in K\}.
\tag{PGR.33}
\]
Every such \(r\) is an original \texttt{GCar}: its rational
isometry \(\sigma_r\) carries \(\Lambda_{u',\rho'}\)
onto \(\Lambda_{u,\rho}\), and its transported flag satisfies
the full original \texttt{PD}, namely \texttt{IsLift},
\texttt{IsAdmissible}, and \texttt{FrameOK}.  For these columns,
the active original terms are exactly
\[
       N_{u,\rho}\backslash
                 \Gamma_{\rho,\rho'}/N_{u',\rho'}.
\tag{PGR.34}
\]
Here ``active'' has its literal meaning: both row lattices match the
columns and the prescribed representative has \(k_r\in K\).
No term is restricted by a nondegeneracy condition on its Gram matrix.
\begin{proof}
The original rational coordinates identify \(r\) with a unique
\(E\)-linear isometry \(\sigma_r\).  The identity
\(r_f a_{u',\rho'}=a_{u,\rho}k_r\), with
\(k_r\widehat L_0=\widehat L_0\), gives equality of the
completed lattices.  The integral-coordinate argument of PGR9 turns
this into equality of the rational lattices themselves.  Thus the
lattice part of \texttt{GCar} holds in both directions.

We spell out the flag conditions, since a bare lattice isometry is
not the source's return carrier.  The supplied row condition for
\(\Lambda_{u',\rho'}\) includes that its negative vector belongs
to that lattice, that its chosen lift is this vector, the prescribed
local convention for the isotropic frame, the integral primitivity
condition on the negative vector in the quotient, and the equality
\(X\Lambda=XS\) with the prescribed different-weighted span.
Applying \(\sigma_r\) preserves all these conditions.  For the
local convention this follows by transporting each localized spanning
identity with the same nonzero coefficients; coefficients required
to be one at primes above \(g\) remain one.  Intersections with
the frame span and its perpendicular, sums with that span, and the
quotient condition defining \(L_{\rm sub}\) are carried to their
corresponding objects because \(\sigma_r\) is an isometry.
Consequently any scalar whose multiple of the transported negative
vector belongs to \(L_{\rm sub}\) is integral exactly when it was
integral before transport.  The different-weighted span transports
with its same coefficient ideals.  These observations prove both
\texttt{IsLift} and every field of \texttt{FrameOK}.

For \texttt{IsAdmissible}, write \(\mathcal F_u\) for the
ordered frame \((e_0,\ldots,e_{r-1},z_u)\).  Above every prime
dividing \(g\), the supplied congruences for \(x_u\) and the
fact that \(h_{u,\rho}\) fixes \(\mathcal F_u\) identify
\(a_{u,\rho}^{-1}\mathcal F_u\) with the reference frame
modulo \(g\widehat L_0\).  The same holds for \(u',\rho'\).
The defining conditions of \(K_2\) say that \(k_r\) acts
upper unitriangularly on the reference \((e,z)\)-frame modulo
\(g\widehat L_0\).  Hence it supplies coefficients, integral
in the finite adeles, with zero entries below the diagonal and one
on the diagonal, for
\[
 a_{u,\rho}^{-1}
       \bigl(\sigma_r\mathcal F_{u'}-\mathcal F_u\xi\bigr)
                 \in g\widehat L_0
                        \quad\hbox{at all primes above }g.
\tag{PGR.35}
\]
The finitely many off-diagonal residue coefficients can be lifted
simultaneously to \(\mathcal O_E\) by Chinese remaindering.
Set the lower entries to zero and the diagonal entries to one
exactly; call the resulting matrix \(\xi\) as well.  Replacing
the finite coefficients by these lifts does not change (PGR.35).

For a rational vector \(v\), the condition
\(a_{u,\rho}^{-1}v\in g\widehat L_0\) at the primes
above \(g\) says precisely that
\[
 v\in(g\Lambda_{u,\rho})_{\mathfrak p}
                         \quad\hbox{for every }\mathfrak p\mid g.
\tag{PGR.36}
\]
This follows by writing coordinates in a rational lattice basis and
localizing; equivalently use the localized congruence interpretation
of ASG11.  The statement (PGR.36), not a stronger global
\(g\Lambda\)-membership assertion, is the source's
\texttt{ModG}.  Thus (PGR.35) with this one integral upper
unitriangular \(\xi\) proves every field of
\texttt{IsAdmissible}.  The entire \texttt{PD} predicate is proved.

Finally, the source relation on two carriers with the same rows is
exactly \(r' = n r n'\), with
\(n\in N\!\operatorname{Lam}(\Lambda_{u,\rho})\) and
\(n'\in N\!\operatorname{Lam}(\Lambda_{u',\rho'})\).
PGR9 identifies these groups, including their unique rational lifts,
with the two groups in (PGR.34).  It also proves that both left and
right multiplication preserve \(k_r\in K\).  Hence activity
does not depend on the chosen carrier representative, and the
quotient is exactly (PGR.34).  Conversely every active original
carrier has its specified rational return in (PGR.33), so nothing
has been lost from this quotient dictionary.
\end{proof}
% END PGR14

% BEGIN PGR15
\paragraph{PGR15. Pointwise support reindexing, with the exact finite phase.}
\label{claim:PGR15}
For every \(T'\), \(\rho\), and \(x\in G_\infty\),
\[
 P_{u'}T'(x,a_{u,\rho})
  =\sum_{\rho'\in I_{u'}}
       \sum_{[r]\in\Gamma_{\rho,\rho'}/N_{u',\rho'}}
         \chi(k_r)^{-1}
           \langle J_{u'}T'(\rho'),W_{u'}(r_\infty^{-1}x)\rangle.
\tag{PGR.37}
\]
The same reindexing holds with absolute values, and both sums are
absolutely convergent.  The right-coset index is also exactly the
original \texttt{ColCl}, originally defined using \(S_{u'}(F)\).
Thus this is the full original \texttt{SupportReindex}, with its
norm-sum assertion as well as its scalar-sum assertion.
\begin{proof}
Invert a left-coset index \(\delta\) in the original period, so
its summand is \(\Phi_{u'}(T')(r^{-1}(x,a_{u,\rho}))\).
The support of this seed depends only on its finite coordinate.
Whenever this coordinate is in the support, the finite part of the
coverage argument in PGR10 supplies a unique column \(\rho'\)
and a rational element of \(S_{u'}(F)\) by which to modify \(r\)
on the right so that
\(r_f^{-1}a_{u,\rho}\in a_{u',\rho'}K\).
Equivalently \(r\in\Gamma_{\rho,\rho'}\).
Left rational invariance of the seed, proved in PGR7, means this
modification does not alter its value at any infinite coordinate.
Cosets outside the support contribute zero and are simply omitted.

Suppose two returns in \(\Gamma_{\rho,\rho'}\) differ on the
right by \(s'\in S_{u'}(F)\).  Their finite relation implies
\(a_{u',\rho'}^{-1}s'_f a_{u',\rho'}\in K\).
Split \(s'=n'h'\) rationally.  Integral Levi extraction, PGR4,
puts the Levi factor in the conjugate compact level of that column.
PGR5 forces its rational Levi part to be one.  PGR9 then says
\(s'\in N_{u',\rho'}\).  The converse follows immediately
from PGR9.  Distinct columns cannot be identified, by uniqueness of
the original Levi column class.  This proves the claimed bijection
of support indices, including the exact \texttt{ColCl} relation.

At a return in (PGR.33) one has the literal finite equality
\(r_f^{-1}a_{u,\rho}=a_{u',\rho'}k_r^{-1}\).
Evaluate the seed using \(n'=1\), \(h'=h_{u',\rho'}\),
and \(\kappa'=k_r^{-1}\).  Formula (PGR.3) gives precisely
the summand on the right of (PGR.37).
Changing \(r\) to \(rn'\), \(n'\in N_{u',\rho'}\),
multiplies \(\chi(k_r)^{-1}\) by
\(\psi_{u',\infty}(n')\), by (PGR.19), and multiplies
the wave-function value by its inverse.  The value is unchanged.
The bijection therefore preserves the scalar summands and their
absolute values.  Pointwise absolute convergence follows from PGR12
(and is also part of the retained \texttt{InHH}).
\end{proof}
% END PGR15

% BEGIN PGR16
\paragraph{PGR16. Orbit unfolding with arbitrary original sections.}
\label{claim:PGR16}
The left action of \(N_{u,\rho}\) on the right-coset set in
(PGR.37) has the active original terms as its orbit set.  If \(r\)
represents a term \(t\), its stabilizer is exactly
\[
 R_r=N_{u,\rho}\cap rN_{u',\rho'}r^{-1},
\tag{PGR.38}
\]
viewed in \(G_\infty\); this is the original
\texttt{ROinf}.  Define
\[
 f_r(x)=\overline{\langle J_uT(\rho),W_u(x)\rangle}
       \chi(k_r)^{-1}
       \langle J_{u'}T'(\rho'),W_{u'}(r^{-1}x)\rangle.
\tag{PGR.39}
\]
Whenever the corresponding absolute slab integral is finite,
\[
 \int_{\Omega_{u,\rho}}
       \sum_{[r]\in\Gamma_{\rho,\rho'}/N_{u',\rho'}}
                 f_r(x)\,d\mu_\infty(x)
   =\sum_{t\ {\rm active}}\int_{F_t} f_{r_t}(x)\,d\mu_\infty(x),
\tag{PGR.40}
\]
where each \(r_t\) is the original representative and each
\(F_t\) is any prescribed strict fundamental domain for its
\(R_{r_t}\).  The integrand on the left and every term on the
right are integrable.  The same equality holds for the nonnegative
sum of their absolute values.  In particular the sum of the absolute
orbital integrals is finite.  This proves the full original
\texttt{OrbitUnfold} in this setting, using its stated absolute
slab premise; PGR13 supplies that premise in the application.
\begin{proof}
The orbit and stabilizer identifications are now algebraic.
Two right cosets lie in the same left orbit if and only if
their returns differ by \(r\mapsto n r n'\), precisely the
relation in PGR14.  A left multiplier fixes \(rN_{u',\rho'}\)
if and only if it belongs to (PGR.38).  Unique rational lifting,
PGR9, identifies this intersection with the source's definition in
infinite coordinates.

We check the scalar covariance that makes the change of domains valid.
The right multiplication calculation in PGR15 gives
\(f_{rn'}=f_r\).  If \(n\in N_{u,\rho}\), (PGR.19)
and (PGR.1) give
\[
                  f_{nr}(x)=f_r(n^{-1}x).
\tag{PGR.41}
\]
Indeed \(\chi(k_{nr})^{-1}=\psi_{u,\infty}(n)
\chi(k_r)^{-1}\), whereas
\(\overline{\langle J_uT(\rho),W_u(n^{-1}x)\rangle}
=\psi_{u,\infty}(n)
\overline{\langle J_uT(\rho),W_u(x)\rangle}\).
These are exactly the same factor.  For \(n\in R_r\), the
first covariance and (PGR.41) show that \(f_r\) is left
\(R_r\)-invariant.  All these identities hold also for absolute
values.

Choose representatives \(n\) of \(N_{u,\rho}/R_r\).
Integration of (PGR.41) over \(\Omega_{u,\rho}\), followed by
\(y=n^{-1}x\), tiles an \(R_r\)-fundamental domain by the
sets \(n^{-1}\Omega_{u,\rho}\).  Apply this first to
\(|f_r|\).  Tonelli, the countability of the rational group, and
the assumed absolute slab bound show
\[
 \sum_{t\ {\rm active}}\int_{F_t}|f_{r_t}|<\infty.
\tag{PGR.42}
\]
Changing any tiled domain to an arbitrary strict domain \(F_t\)
is legitimate by the same measurable intersection partition as in
PGR10.  Only left Haar invariance is used.  Finite (PGR.42) permits
the complex-valued Fubini interchange and proves (PGR.40).

One can choose the representative of each orbit to be the prescribed
representative of the original term: it is active by PGR14.
For clarity, the operator version of its choice independence has the
same cancellation.  Under \(r\mapsto n r n'\), changing
variable \(x=ny\) in the rank-one orbital integral multiplies
that operator by
\(\psi_{u,\infty}(n)\psi_{u',\infty}(n')\).
Its finite factor \(\chi(k_r)\) acquires the inverse multiplier,
by (PGR.19).  Thus the phase-weighted operator is unchanged.
The supplied \texttt{AmpWellDef} and \texttt{OOFundIndep}
assert precisely the permitted independence for every original
carrier and every original strict domain, including the prescribed
\texttt{FundOf}.  We retain these original premises throughout;
our tiling does not replace either by a special selection of domains
or by a different quotient measure.
\end{proof}
% END PGR16

% BEGIN PGR17
\paragraph{PGR17. Exact term evaluation and the full entry identity.}
\label{claim:PGR17}
For an original term \(t\), keep its prescribed representative
\(r_t\) and its prescribed domain \(F_t=\texttt{FundOf}(t)\).
Its operator is exactly
\[
 O_t=\int_{F_t}
      \bigl[\xi\longmapsto
       \langle W_{u'}(r_t^{-1}x),\xi\rangle W_u(x)\bigr]
                                         \,d\mu_\infty(x).
\tag{PGR.43}
\]
Let \(B_{t,\rho\rho'}\) be zero if the term is inactive for
these columns, and otherwise set
\[
 B_{t,\rho\rho'}=s_{u,u'}v_K\chi(k_{r_t})O_t.
\tag{PGR.44}
\]
These are the literal original \texttt{blk} operators.  Under
the retained \texttt{AbsConvT}, every block series is absolutely
convergent in operator norm, and
\[
 s_{u,u'}\langle P_uT,P_{u'}T'\rangle
   =\sum_{\rho,\rho'}
      \left\langle
       \left(\sum_t B_{t,\rho\rho'}\right)J_{u'}T'(\rho'),
                         J_uT(\rho)\right\rangle
                   \quad\hbox{for all }T,T'.
\tag{PGR.45}
\]
This is the full original \texttt{EntryUnfold} on this setting.
\begin{proof}
The supplied \texttt{OOIntAll} makes (PGR.43) a Bochner integral
of the original rank-one operator, in a finite-dimensional normed
operator space.  For fixed vectors \(v=J_uT(\rho)\) and
\(v'=J_{u'}T'(\rho')\), its matrix coefficient is
\[
 \left\langle O_t v',v\right\rangle
  =\int_{F_t}
       \langle v',W_{u'}(r_t^{-1}x)\rangle
       \langle W_u(x),v\rangle\,d\mu_\infty(x).
\tag{PGR.46}
\]
Indeed evaluation followed by this inner product is a continuous
real-linear functional of the operator, so commutes with the
Bochner integral.  Conjugate linearity in the first variable gives
\(\langle\langle W',v'\rangle W,v\rangle
=\langle v',W'\rangle\langle W,v\rangle\), proving (PGR.46).
In particular an active term satisfies
\[
 s_{u,u'}v_K\int_{F_t}f_{r_t}(x)\,d\mu_\infty(x)
       =\langle B_{t,\rho\rho'}v',v\rangle.
\tag{PGR.47}
\]
Here the scalar in (PGR.44) is conjugated on the left side of the
inner product: \(\overline{\chi(k_{r_t})}=\chi(k_{r_t})^{-1}\).
This accounts for the entire phase in (PGR.39).  The factor
\(s_{u,u'}v_K\) is real, so is not altered.  Equation (PGR.47)
is the exact original \texttt{TermVal}.

Insert the support identity PGR15 into (PGR.32), then apply the
orbit unfolding PGR16 separately for each pair of columns.  PGR13
supplies its absolute hypothesis for every such pair.  The resulting
scalar series is absolutely convergent by (PGR.42).  Formula
(PGR.47) identifies every term, including zero terms, with its
prescribed block coefficient.  Thus the left side of (PGR.45) equals
\(\sum_{\rho,\rho'}\sum_t
\langle B_{t,\rho\rho'}v',v\rangle\).

We now justify the operator sums in (PGR.45); this is a different
convergence issue from the absolute unfolding just completed.  Write
\(A_t\) for the original \texttt{aTerm}, the finite block
operator on \(\bigoplus_{\rho'}V_{u'}\) with blocks
\(B_{t,\rho\rho'}\), and \(\operatorname{op}(t)\) for its
original Hilbert--Schmidt encoding.  The first clause of
\texttt{AbsConvT} supplies
\(\sum_t\|\operatorname{op}(t)\|<\infty\).
In these fixed finite-dimensional spaces there is a finite constant
\(C\), independent of \(t\), with
\[
 \|B_{t,\rho\rho'}\|\leq\|A_t\|
                        \leq C\|\operatorname{op}(t)\|.
\tag{PGR.48}
\]
The first inequality composes with a norm-one coordinate inclusion
and projection (and is trivial for a zero coordinate space).  For
the second, choose orthonormal bases, recover the finitely many
matrix entries from the Hilbert--Schmidt encoding, and bound the
operator norm by the sum of their absolute values; Cauchy--Schwarz
supplies a fixed dimension-dependent constant.  This also proves
the assertion when an operator space is zero, without requiring a
positive dimension.  Hence both \(\sum_tA_t\) and every block
series converge absolutely.  Continuous evaluation, the inner
product, and the finite sums over columns commute with these
convergent sums.  This proves (PGR.45) with exactly its displayed
operator series.
\end{proof}
% END PGR17

% BEGIN PGR18
\paragraph{PGR18. The full original universal \texttt{KExpansionC}.}
\label{claim:PGR18}
For every original \(D\), every supplied \(X\), every
\(Y:\texttt{Full}(X)\), and every ordered pair of words
\(u,u'\), the original assumptions
\[
 \begin{gathered}
 \texttt{FoldOK}(u),\quad\texttt{FoldOK}(u'),\quad
 \texttt{InHH}(u),\quad\texttt{InHH}(u'),\\
 \texttt{FundExists},\quad\texttt{OOIntAll},\quad
 \texttt{OOFundIndep},\quad\texttt{AmpWellDef},\quad
 \texttt{AbsConvT}
 \end{gathered}
\tag{PGR.49}
\]
imply the original \texttt{KExpansion} for
\(X.Fu(u),X.Fu(u'),Y.cs(u),Y.cs(u'),Y.FN(u),Y.FN(u')\),
the given character \(Y.hK\), and the representatives
\(Y.Dt(u),Y.Dt(u')\).  More explicitly, there exists a
continuous complex-linear operator \(B_{u,u'}\) such that
\[
 \operatorname{HasSum}_{t}\,A_t=B_{u,u'},\qquad
 s_{u,u'}\langle P_uT,P_{u'}T'\rangle
       =\langle B_{u,u'}J_{u'}T',J_uT\rangle
                                  \quad\hbox{for every }T,T'.
\tag{PGR.50}
\]
Every object in (PGR.50) is the original supplied object.  The two
words may coincide; no lower bound on their weights or their gap
has been added.
\begin{proof}
Define \(B_{u,u'}=\sum_tA_t\).  Estimate (PGR.48) gives an
absolutely convergent series in the complete finite-dimensional
space of continuous operators.  Equivalently, its net of finite
partial sums converges, which is precisely the \texttt{HasSum}
assertion.  Each block projection is continuous, so the blocks of
\(B_{u,u'}\) are the block sums in (PGR.45).  The inner product
of finite direct sums now turns (PGR.45) into the second identity
of (PGR.50).

We record explicitly how the scope matches the original interface.
The arbitrary \texttt{Full} supplies both column systems, both
row systems, both rational coset selections, and both unipotent
domains.  PGR1--PGR10 prove the required finite phase, splitting,
trivial column stabilizer, and support-cell statements for this
arbitrary setting.  The two \texttt{FoldOK} inputs give the common
averaging set and the absolute bound in PGR12--PGR13, with their
original quantifiers over all column vectors and all domains.
The two \texttt{InHH} inputs identify the periods in the actual
automorphic Hilbert space.  The four orbital inputs in (PGR.49)
are retained at their full original scopes, including every
carrier pair and every allowed section.  PGR14--PGR17 compare to
their prescribed representatives, domains and integrands.
Finally \texttt{AbsConvT} supplies the operator-norm convergence
used here.  Its other clauses remain supplied rather than being
silently weakened in the target statement.

The normalization is exactly (PGR.5), using the given
\(\mu_N,\mu_H,F_N\) for each word.  Those measures and
domains were never replaced.  The unnormalized scalar identity
was proved first, so even if a real-valued mass in the literal
interface is zero, the total inverse convention simply makes both
normalized sides zero.  No proof step divides by that mass.
Only the finite level volume was divided out in PGR13, and its
strict positivity and finiteness hold for every supplied setting.
This verifies all the quantifiers in \texttt{KExpansionC}.
\end{proof}
% END PGR18

\subsection{Transport to the actual period operator}

% BEGIN PGR19
\paragraph{PGR19. Linearity, the antilinear Gram transport, and projected periods.}
\label{claim:PGR19}
For every word and every permitted \(D_u\), the actual
\texttt{PuLinear} assertion holds: under its original
\texttt{InHH} premise, \(P_u\) is complex linear.
Give \(C_u\) its original \texttt{ColT} inner product from
PGR6, and regard \(P_u\) as a map into the actual ambient
quotient \(L^2\) space \(\mathcal H\).  Put
\[
 \mathcal A_u=(\sqrt{c_u})^{-1}P_u.
\tag{PGR.51}
\]
Under the hypotheses of PGR18, the exact Gram transport is
\[
 B_{u,u'}=J_u\,\mathcal A_u^*\mathcal A_{u'}\,J_{u'}^{-1}.
\tag{PGR.52}
\]
The two occurrences of antilinearity in (PGR.52) are essential.
For any subspaces \(E_u\subset C_u\), \(E_{u'}\subset C_{u'}\),
with inclusions \(\iota_u,\iota_{u'}\), and any orthogonal
projection \(Q\) on a closed subspace of \(\mathcal H\),
the Gram operator of the projected periods is exactly
\[
 \begin{split}
 (Q\mathcal A_u\iota_u)^*(Q\mathcal A_{u'}\iota_{u'})
   &=\iota_u^*\mathcal A_u^*\mathcal A_{u'}\iota_{u'}\\
   &\quad-\iota_u^*\mathcal A_u^*(1-Q)
                              \mathcal A_{u'}\iota_{u'}.
 \end{split}
\tag{PGR.53}
\]
Thus a subsequent target projection may be suppressed only after
its displayed correction term has been proved zero.
\begin{proof}
For any fixed finite support decomposition, (PGR.3) is linear in
\(T\), and PGR7 says that its value is independent of that
decomposition.  Off support it is zero.  Hence the actual scalar
seed is linear in \(T\) at every point.  Under \texttt{InHH},
all the scalar Poincar\'e sums converge absolutely.  One can
therefore add them and multiply them by arbitrary complex scalars
term by term.  Their \(L^2\) representatives satisfy the same
identities because \texttt{InHH} supplies their \(L^2\)
membership.  This proves the full implication defining
\texttt{PuLinear}; no folding or gap hypothesis is needed for
this linearity argument.  Its range in the prescribed automorphic
subspace is the range assertion already supplied by \texttt{InHH}.

PGR6 makes \(C_u\) a finite-dimensional Hilbert space with
\(J_u\) an antilinear isometry onto the complete finite direct
sum of compact-type spaces.  Thus the linear map \(P_u:C_u\to
\mathcal H\) is continuous, and its Hilbert adjoint exists.
We use the ambient \(\mathcal H\) here; no unstated completeness
of a proper automorphic subspace is required.

Each real-valued Haar mass in (PGR.5) is nonnegative, as are
\(v_K\) and the natural-number dimension.  Total division
therefore gives \(c_u\geq0\).  For two such numbers the
square-root product and inverse identities, also when either is
zero, give
\[
 \langle\mathcal A_uT,\mathcal A_{u'}T'\rangle
             =s_{u,u'}\langle P_uT,P_{u'}T'\rangle.
\tag{PGR.54}
\]
Set \(C=\mathcal A_u^*\mathcal A_{u'}\).  The adjoint
identity, the original \texttt{ColT} inner product, and PGR18
give, for all \(T,T'\),
\[
 \langle J_u C T',J_uT\rangle
   =\langle T,CT'\rangle
   =\langle\mathcal A_uT,\mathcal A_{u'}T'\rangle
   =\langle B_{u,u'}J_{u'}T',J_uT\rangle.
\]
Surjectivity of \(J_u\) and nondegeneracy of the inner product
identify the first vectors.  Surjectivity of \(J_{u'}\)
proves (PGR.52).  In particular, this is not the formula obtained
by treating the Riesz map as complex linear.

For (PGR.53), use \(Q^*=Q=Q^2\) and the rule for the adjoint
of a composite.  The left side is
\(\iota_u^*\mathcal A_u^*Q\mathcal A_{u'}\iota_{u'}\).
Substitute \(Q=1-(1-Q)\).  This proves (PGR.53) with no
additional range assertion.  If \(Q\) fixes the relevant period
images, its correction vanishes; otherwise the correction must be
kept in any subsequent comparison.
\end{proof}
% END PGR19

% BEGIN PGR20
\paragraph{PGR20. The full original \texttt{KTie} consumer on the supplied finite datum.}
\label{claim:PGR20}
Retain an arbitrary original \(h_F:\texttt{FinX}(Y)\), and let
\(s_Y=\texttt{sOf}(h_F)\) be its actual constructed spectral
datum.  For every pair of words, the original two
\texttt{DimTie} identities, \texttt{KExpansion}, and
\texttt{AbsConvT} imply
\[
     \texttt{KExpansion}(s_Y,u,u')
                     \quad\hbox{and}\quad
     \texttt{KAbsConv}(s_Y,u,u').
\tag{PGR.55}
\]
This is the full original \texttt{KTie}; in particular PGR18
now supplies its analytic expansion input at universal scope.
The fields of \texttt{FinX} remain the original ones: the two
column bijectivity fields, \texttt{InHH} and period linearity
for all words, finite dimension of the actual automorphic space,
finite fibers of the map to lattice terms, pair independence,
the real-value field, the coefficient-scale supplier, and volume
positivity.  No existence assertion for \texttt{FinX}, and no
isomorphism from a different spectral datum, is inferred here.
\begin{proof}
We give the transport and regrouping rather than treating
\texttt{KTie} as a consequence of notation.  The column spaces
of \(s_Y\) are the original \texttt{ColT} spaces, with the
same inner products as PGR6.  Its period maps are precisely
\(P_u\), now viewed in the supplied finite-dimensional actual
automorphic space.  Let \(B\) be the operator provided by the
assumed \texttt{KExpansion}, and let
\[
       \mathcal T(A)=J_u^{-1} A J_{u'}
\tag{PGR.56}
\]
for a continuous complex-linear operator between the two column
coordinate sums.  This is again a complex-linear operator on the
original \texttt{ColT} spaces.  As a function of \(A\),
\(\mathcal T\) is additive and real linear; it conjugates
complex scalar multiplication.  Both \(J\)'s are isometries
and surjective, so
\(\|\mathcal T(A)\|=\|A\|\).
In particular it is a continuous additive map and preserves
convergent sums.

The original normalizing constant of \(s_Y\) is
\(Y.\operatorname{volN}(u)Y.\operatorname{volH}(u)/
\operatorname{dimPi}(u)\).  Its numerator is exactly
\(v_K\mu_{N,u}(F_{N,u})\mu_{H,u}(H_{u,\infty})\),
with the same real-valued mass convention as before.  The
\texttt{DimTie} assumption identifies its denominator with
\(\dim_{\mathbb C}V_u\).  Thus it equals \(c_u\) exactly,
and likewise for \(u'\).  Nonnegativity identifies the source's
real power \((c_uc_{u'})^{-1/2}\) with
\((\sqrt{c_uc_{u'}})^{-1}\), including the zero convention.
Testing against arbitrary column vectors as in PGR19 now proves
\[
                 K_{s_Y}(u,u')=\mathcal T(B).
\tag{PGR.57}
\]
The finite-dimensionality supplied by \(h_F\) makes all
adjoints in this original definition legitimate.  This calculation
uses no identification of \(s_Y\) with the initial datum \(D.s\).

Let \(\mathfrak T\) be the original two-row orbital term set
and \(\mathcal R\) the original lattice-flag term set.  The
map \(q:\mathfrak T\to\mathcal R\) is the source's
\texttt{toTerm}; it transports the primed flag and forgets the
second row presentation.  For each \(t\in\mathcal R\),
the actual \(s_Y\)-term operator is defined to be
\[
             A^R_t=\sum_{\tau:\ q(\tau)=t}\mathcal T(A_\tau).
\tag{PGR.58}
\]
This is the literal finite sum, since \(h_F.\texttt{fibFin}\)
supplies finiteness for every fiber.  The absolute convergence
clause and (PGR.48) give
\(\sum_{\tau\in\mathfrak T}\|A_\tau\|<\infty\),
and hence the same assertion for \(\mathcal T(A_\tau)\).
Applying the continuous additive map \(\mathcal T\) to the
original \texttt{HasSum} gives sum \(\mathcal T(B)\).
Partition this absolutely convergent series into its fibers under
\(q\).  The triangle inequality gives
\[
 \sum_{t\in\mathcal R}\|A^R_t\|
       \leq\sum_{\tau\in\mathfrak T}\|\mathcal T(A_\tau)\|<\infty,
 \qquad
       \operatorname{HasSum}_{t\in\mathcal R}A^R_t=\mathcal T(B).
\tag{PGR.59}
\]
Empty fibers contribute zero.  One may prove the regrouping directly
by choosing a finite set of \(\tau\)'s with arbitrarily small
norm tail; every finite regrouping containing their finitely many
images has the same bound on its remaining tail.  Thus the assertion
is independent of any enumeration or choice of fiber representatives.
Equations (PGR.57)--(PGR.59) are exactly the two conclusions in
(PGR.55).  The additional fields of \(\texttt{FinX}\) remain
supplied for the construction of \(s_Y\), even where this
particular regrouping uses only a subset of them.
\end{proof}
% END PGR20

\subsection{What this supplies, and what remains}

PGR18 supplies the full original universal \texttt{KExpansionC},
and PGR20 supplies its original \texttt{KTie} transport with the
actual common data.  The proof has also reconstructed the finite
phase, integral Levi extraction, trivial column stabilizers, the
original flag return predicate, absolute unfolding, and the exact
matrix-coefficient evaluation.  Each of these intermediate arguments
is retained above so that the identity can be read without following
the source's chain of wrappers.

This is an identity under the original analytic inputs.  It does not
produce \texttt{Full}, all-word folding and automorphic-range
suppliers, \texttt{AbsConvC}, or the remaining fields needed to
construct \texttt{FinX}.  It does not construct a common-family
\texttt{SIso}, remove the projected-period correction in (PGR.53),
or establish an upper bound for the resulting Gram operator.
The exact archimedean Jacobian and limiting-integral producers
\texttt{JacGauge} and \texttt{RhsWLim}, the independent product
measure statement \texttt{ProdExact}, and the remaining mass,
counting and spectral comparisons stay separate.  No whole global
spectral obligation is closed by this chapter alone.
